\src{183}
Since (5.3 bis) is an isomorphism,\footnote{The printed text gives a second
formula ``(5.3)'' on page 182 and refers to it here; the identifier
``(5.3 bis)'' distinguishes this formula from the first (5.3).} so is (5.4),
and hence so is (5.5). Since $s_{(x)}$ is a basis of
$\operatorname{Ext}_{\Osh_X}^{n}(X;k_x,\Omega_X^{n})$ by (4.2), it follows that
its image $\varepsilon_X$ is a basis of $H^{n}(X,\Omega_X^{n})$.

It remains to prove assertion (D), which will follow in a purely formal way
from the elementary facts summarized in the following lemmas. Here $X$ is
assumed to be a closed subset (singular or not) of projective space $P$ of
dimension $r$. We use $\Osh_P(m)$ for the sheaf on $P$ denoted $\Osh(m)$ in
[3], and $\Osh_X(m)$ for the analogous sheaf on $X$.

\result{LEMMA}{2} Assertion (D) is true if $X=P$ and $F=\Osh_P(m)$.

This lemma is verified by a direct calculation. The explicit calculation of
$H^{i}(P,\Osh_P(m))$ is given in [3], but it can be carried out more
elementarily. The computation of the cup product
\[
  H^{i}(P,\Osh_P(m))\times H^{j}(P,\Osh_P(m'))
  \longrightarrow H^{i+j}(P,\Osh_P(m+m'))
\]
needed to make the pairing (5.1) explicit presents no
difficulty.\footnote{The printed text has $\Osh_P(r+r')$; the two factors
$\Osh_P(m)$ and $\Osh_P(m')$ require $\Osh_P(m+m')$.}

\result{LEMMA}{3} Every coherent algebraic sheaf $F$ on $X$ is isomorphic to a
quotient sheaf of a sheaf $\Osh_X(-m)^{k}$, where $m$ may be taken as large as
desired.

This follows from the fact that $F\otimes\Osh_X(m)$ is ``generated by its
sections'' for $m$ large; cf. [3].

\result{LEMMA}{4} Let $i>0$. Then
$H^{r-i}(P,\Osh_P(-m))=0$ for $m$ large; and for every coherent algebraic sheaf
$B$ on $X$, one has
\[
  \operatorname{Ext}_{\Osh_X}^{i}(X;\Osh_X(-m),B)=0
\]
for $m$ large.

The first assertion follows from the explicit calculations mentioned above.
For the second, note that there is an isomorphism
\[
  \operatorname{Ext}_{\Osh_X}^{i}(X;\Osh_X(-m),B)
  =H^{i}(X,B\otimes\Osh_X(m))
\]
(Proposition 1, Corollary 1), whence the conclusion by a well-known result of
[3]. Combining the two preceding lemmas, we obtain the

\result{COROLLARY}{} Let $i>0$. Then the functor
$F\longmapsto H^{r-i}(P,F)$ on the category of coherent algebraic sheaves on
$P$ is coeffaceable, and so is the functor
$\operatorname{Ext}_{\Osh_X}^{i}(X;F,B)$ on the category of coherent algebraic
sheaves on $X$.

